% ============================================================
% phu-luc-b-ma-nguon.tex — Phụ lục mã nguồn: mã nguồn, lệnh chạy lại, phân công đọc hiểu code, (Giang)
% ============================================================
\chapter{MÃ NGUỒN VÀ DỮ LIỆU}
\label{app:code}

\section*{Cấu trúc thư mục}

Mã nguồn viết bằng Python 3.9 với NumPy, SciPy, pandas và scikit-learn (phiên bản ghim trong
\texttt{requirements.txt}). FCM và các phiên bản PSO được cài bằng NumPy. scikit-learn chỉ dùng cho Silhouette và AUC, còn
scikit-fuzzy chỉ dùng trong phép thử để đối chiếu FCM. Bảng~\ref{tab:pl-ma-nguon} liệt kê các tệp dùng cho kết quả trong
bài. Thư mục \texttt{code/} còn chứa các phương pháp của phiên bản trước (thuật toán GWO, các cách ghép khác), không trình bày ở đây.

{\footnotesize
\begin{longtable}{@{}>{\raggedright\arraybackslash}p{4.4cm}>{\raggedright\arraybackslash}p{10.4cm}@{}}
\caption{Các tệp mã nguồn dùng trong bài}\label{tab:pl-ma-nguon}\\
\toprule
\textbf{Tệp} & \textbf{Nội dung} \\
\midrule
\endfirsthead
\toprule
\textbf{Tệp} & \textbf{Nội dung} \\
\midrule
\endhead
\texttt{fcm.py} & FCM (mục~\ref{sec:fcm}), khởi tạo U ngẫu nhiên hoặc k-means++ \\
\texttt{pso\_co\_so.py} & PSO có vận tốc dùng chung cho PSO gốc, PSO-W, PSO-S, PSO-LOG (mục~\ref{sec:pso-goc}) \\
\texttt{thuat\_toan\_a.py}, \texttt{thuat\_toan\_b.py} & Lịch quán tính S và LOG (mục~\ref{sec:pso-s-log}); QPSO (mục~\ref{sec:qpso}) \\
\texttt{cach\_ket\_hop.py} & Bài toán trọng số (fitness Silhouette trên không gian gốc, phạt cụm nhỏ), phân hoạch theo trọng số, WFCM giải tích (mục~\ref{sec:wfcm}, \ref{sec:pso-trong-so-pp}) \\
\texttt{pso\_trong\_so.py} & PSO tìm trọng số thuộc tính với 5 phiên bản, tìm ngẫu nhiên cùng ngân sách (mục~\ref{sec:pso-trong-so-pp}) \\
\texttt{pso\_khoi\_tao.py} & PSO tìm tâm khởi tạo cho FCM, 5 phiên bản, đối chứng điểm ngẫu nhiên (mục~\ref{sec:pso-khoi-tao}) \\
\texttt{giao\_dien.py} & Giao diện chung \texttt{fit(X, c, m, rng, max\_fe)}, FCM k-means++, FCM tinh chỉnh từ bộ tâm \\
\texttt{chi\_so\_danh\_gia.py}, \texttt{danh\_gia.py} & PC, PE, XB, Silhouette, AUC, chênh tỉ lệ vỡ nợ, đánh giá một lượt \\
\texttt{tien\_xu\_ly.py}, \texttt{mo\_ta\_du\_lieu.py} & Nạp, tiền xử lý dữ liệu (mục~\ref{sec:tien-xu-ly}), 2 nửa bộ Việt Nam theo doanh nghiệp, hồ sơ chất lượng các bộ \\
\texttt{chay\_e5.py}, \texttt{phan\_tich\_e5.py}, \texttt{chay\_e7.py} & TN1 (mã E5, E7; PSO tìm tâm, cấu hình nhiều cực tiểu, quỹ đạo gbest) \\
\texttt{chay\_thuc\_nghiem.py}, \texttt{phan\_tich.py} & TN2 phần chẩn đoán 7 bộ (mã E3), profile cụm (mã E4) và phân tích tương ứng \\
\texttt{chay\_e6.py} & TN2 phần quét $m$, $c$ (mã E6): bảng, macro \texttt{results/so-lieu-e6.tex} \\
\texttt{chay\_e8.py} & TN3 (mã E8; PSO-W tìm trọng số so với 3 đối chứng): bảng, hình, macro \texttt{results/so-lieu-e8.tex} \\
\texttt{chay\_e8\_kmeans.py} & TN3: K-means không trọng số làm mốc tham chiếu, Holm riêng; macro \texttt{results/so-lieu-e8-kmeans.tex} \\
\texttt{chay\_e9.py} & TN4 (mã E9; 5 phiên bản PSO tìm trọng số, 2 ngân sách): bảng, hình, macro \texttt{results/so-lieu-e9.tex} \\
\texttt{chay\_e10.py} & TN5 (mã E10; trọng số học trên mẫu nguồn áp lên mẫu đích): bảng, macro \texttt{results/so-lieu-e10.tex} \\
\texttt{kiem\_thu\_*.py} & 9 tệp phép thử, tổng 97 phép thử; mẫu hồi quy tiền xử lý \texttt{mau\_kiem\_thu/hoi-quy-X.npz} sinh bằng \texttt{tao\_mau\_hoi\_quy.py} \\
\bottomrule
\end{longtable}
}

\section*{Khả dụng dữ liệu và mã nguồn}

Bảng~\ref{tab:pl-du-lieu} liệt kê nguồn và giấy phép của các bộ dữ liệu dùng trong bài. Tệp thô nằm trong
\texttt{data/raw/}, mỗi bộ một thư mục; \texttt{data/README.md} ghi đường dẫn tải, ngày tải và mã SHA-256 của từng tệp
nén. Mã không sửa tệp thô; mọi bước tiền xử lý chạy trong bộ nhớ. Mẫu con 5.000 quan sát được lấy phân tầng theo nhãn với
seed cố định: mẫu thiết kế Taiwan Default dùng seed 0, mẫu kiểm chứng dùng seed 1 trên phần còn lại; 2 bộ phá sản dùng
seed 0 (\texttt{tien\_xu\_ly.py}). Bộ Việt Nam được tác giả công bố kèm bài báo; kho dữ liệu không ghi giấy phép riêng nên
nhóm chỉ dùng để nghiên cứu và không phân phối lại. Mã nguồn, các tệp CSV kết quả và bảng, hình, macro sinh tự động nằm
trong thư mục tiểu luận theo cấu trúc ở Bảng~\ref{tab:pl-ma-nguon}.

{\footnotesize
\begin{longtable}{@{}>{\raggedright\arraybackslash}p{3.6cm}>{\raggedright\arraybackslash}p{4.6cm}>{\raggedright\arraybackslash}p{2.4cm}>{\raggedright\arraybackslash}p{3.8cm}@{}}
\caption{Nguồn và giấy phép dữ liệu}\label{tab:pl-du-lieu}\\
\toprule
\textbf{Bộ dữ liệu} & \textbf{Nguồn} & \textbf{Giấy phép} & \textbf{Dùng ở} \\
\midrule
\endfirsthead
\toprule
\textbf{Bộ dữ liệu} & \textbf{Nguồn} & \textbf{Giấy phép} & \textbf{Dùng ở} \\
\midrule
\endhead
Việt Nam (khó khăn tài chính) & Trần và cộng sự~\parencite{tran2022distress} & không ghi riêng & TN1--TN5, profile cụm \\
Taiwan Default & UCI id 350~\parencite{uciTaiwanDefault2009} & CC BY 4.0 & TN1--TN5, profile cụm \\
South German Credit & Grömping~\parencite{gromping2019sgc}, UCI id 573 & CC BY 4.0 & TN2, TN3 \\
German Credit (bản số) & UCI id 144~\parencite{uciGermanStatlog1994} & CC BY 4.0 & TN2, TN3 \\
Australian Credit & UCI id 143~\parencite{uciAustralian1987} & CC BY 4.0 & TN2, TN3 \\
Polish Bankruptcy & UCI id 365~\parencite{uciPolish2016} & CC BY 4.0 & TN2, TN3 \\
Taiwan Bankruptcy & UCI id 572~\parencite{uciTaiwanBankruptcy2020} & CC BY 4.0 & TN2, TN3 \\
\bottomrule
\end{longtable}
}

\section*{Chạy lại thực nghiệm}

Các lệnh dưới đây chạy trong thư mục \texttt{code/}.
\begin{Verbatim}[fontsize=\footnotesize,frame=single]
pip install -r ../requirements.txt
for f in kiem_thu_*.py; do python3 $f; done        # 97 phép thử
python3 chay_e5.py --tien-trinh 8                  # TN1
python3 phan_tich_e5.py                            # bảng, hình, so-lieu-e5.tex
python3 chay_e7.py --tien-trinh 12                 # TN1, nhiều cực tiểu (gồm phân tích)
python3 chay_thuc_nghiem.py --thi-nghiem E3        # TN2, chẩn đoán 7 bộ
python3 chay_e6.py --tien-trinh 12                 # TN2, quét m, c (gồm phân tích)
python3 chay_e8.py --tien-trinh 16                 # TN3 (gồm phân tích)
python3 chay_e8_kmeans.py --tien-trinh 8           # TN3, mốc K-means (gồm phân tích)
python3 chay_e9.py --tien-trinh 12                 # TN4 (gồm phân tích)
python3 chay_e10.py --tien-trinh 12                # TN5 (gồm phân tích)
python3 chay_thuc_nghiem.py --thi-nghiem E4        # profile cụm
python3 mo_ta_du_lieu.py ../results/mo-ta-du-lieu.csv  # mô tả dữ liệu
python3 phan_tich.py                               # bảng TN2 (E3), profile cụm, so-lieu.tex
cd .. && ./compile.sh                              # dựng build/main.pdf
\end{Verbatim}
Trên máy của nhóm (Intel Xeon E-2278G, 16 luồng, một luồng BLAS mỗi tiến trình), phần chẩn đoán của TN2 và profile cụm
mất dưới một phút, phần quét $m$, $c$ dưới 10 phút; TN1 khoảng 25 phút; TN3 khoảng 2 giờ 20 phút; TN4 và TN5 lâu nhất vì mỗi
lần tính fitness là một lần chạy FCM và một lần tính Silhouette (thời gian ghi trong \texttt{results/log-chay-e*.txt}). Seed cố
định là 1000--1029 nên chạy lại cho cùng kết quả, trừ cột thời gian. Các trình chạy bỏ qua tác vụ đã có trong CSV, nên khi bị ngắt chỉ cần chạy
lại cùng lệnh. Tệp \texttt{results/kiem-tra.txt} lưu các kiểm tra tự động trên kết quả chẩn đoán (TN2), gồm chẩn đoán tỉ lệ 39/100 nêu ở
mục~\ref{sec:vi-sao}.

\section*{Phân công đọc hiểu mã nguồn}

Mỗi thành viên nhận một phần mã nguồn để đọc hiểu, chạy phép thử và trình bày
khi vấn đáp (Bảng~\ref{tab:pl-phan-cong}).

\begin{table}[H]
\centering
\caption{Phân công đọc hiểu mã nguồn}\label{tab:pl-phan-cong}
\small
\begin{tabular}{@{}l>{\raggedright\arraybackslash}p{7.6cm}>{\raggedright\arraybackslash}p{3.6cm}@{}}
\toprule
\textbf{Người} & \textbf{Phần mã đọc hiểu} & \textbf{Mục trình bày} \\
\midrule
Bình & FCM, WFCM giải tích, quy trình đề xuất, kiểm chứng trên mẫu độc lập & \ref{sec:fcm}, \ref{sec:wfcm}, \ref{sec:de-xuat}, \ref{sec:kiem-chung} \\
Giang & PSO gốc, PSO-W, khung PSO tìm trọng số, tiền xử lý dữ liệu, lý do chuyển hướng, TN3 & \ref{sec:pso-goc}, \ref{sec:pso-w}, \ref{sec:pso-trong-so-pp}, \ref{sec:du-lieu}, \ref{sec:chuyen-huong}, \ref{sec:pso-trong-so} \\
Hà & PSO-S, PSO-LOG, chỉ số đánh giá, kiểm định Wilcoxon, Holm, TN4 phần so phiên bản, profile cụm, đánh giá chung & \ref{sec:pso-s-log}, \ref{sec:chi-so}, \ref{sec:kiem-dinh}, \ref{sec:phien-ban-trong-so}, \ref{sec:dien-giai}, \ref{sec:danh-gia-chung} \\
Hải & QPSO, PSO tìm tâm (TN1), chẩn đoán số cực tiểu (TN2), quỹ đạo gbest, chi phí & \ref{sec:qpso}, \ref{sec:pso-khoi-tao}, \ref{sec:e5-dat-j}, \ref{sec:vi-sao} \\
\bottomrule
\end{tabular}
\end{table}
